\documentclass[11pt,a4paper]{report}

\input{ccsit-preamble}
%% =====================================================================
%%                 CCSIT / KFU MACROS (EDIT THESE)
%% =====================================================================
\newcommand{\ksa}{Kingdom of Saudi Arabia}
\newcommand{\moe}{Ministry of Education}
\newcommand{\kfu}{King Faisal University}
\newcommand{\ccsitCollege}{CCSIT, King Faisal University}

\newcommand{\ccsitFull}{College of Computer Sciences \& Information Technology}

\newcommand{\gpTitle}{AI Based Research Documents Repository Using Local Language Models}
\newcommand{\gpTitleShort}{Multimodal Document Repository} % used in header (keep short)
\newcommand{\gpDegree}{Bachelor of Science in Computer Science}
\newcommand{\gpDepartment}{Department of CS}

\newcommand{\gpSemester}{Second Semester}
\newcommand{\gpAcademicYear}{2025--2026}

%% Automatic Month Year based on compile date
\newcommand{\gpSubmissionMonth}{%
\ifcase\month\or
January\or February\or March\or April\or May\or June\or
July\or August\or September\or October\or November\or December%
\fi\space \number\year}

%% Students. For a team of fewer than four, do NOT delete or blank these definitions. Delete the matching lines from the "Prepared by" table in
%% cover.tex and from the signature table in main.tex instead.
\newcommand{\gpStudentA}{Raihanah Nezar Al Ali}
\newcommand{\gpStudentAID}{223007420}
\newcommand{\gpStudentB}{Maryam Ali Alkathim}
\newcommand{\gpStudentBID}{223021603}
\newcommand{\gpStudentC}{Fadak Hani Alibrahim}
\newcommand{\gpStudentCID}{223022618}

%% Supervision / Committee
\newcommand{\gpSupervisor}{Dr. Nora Abdullah Alkhaldi}
\newcommand{\gpCommitteeA}{Dr. Hala Hamdoun}
%% ----------------------------- Acknowledgment Text -----------------------------
\newcommand{\gpAcknowledgmentText}{
We would like to express our sincere gratitude to our supervisor, \textbf{\gpSupervisor}, for the continuous
guidance, valuable feedback, and support throughout the development of this graduation project. We
also extend our appreciation to the committee members and all faculty and staff in \ccsitFull\ for
their assistance and encouragement.
}
\input{undertaking}

%% =====================================================================
%%                           DOCUMENT
%% =====================================================================
\begin{document}

%% ----------------------------- Title Page -----------------------------
%% =====================================================================
%% cover.tex  --  the title page
%% ---------------------------------------------------------------------
%% Built entirely from the macros in ccsit-macros.tex. To change what
%% appears here, edit those macros.
%%
%% Two exceptions you may edit in this file:
%%   - teams of fewer than four students: delete the unused lines from
%%     the "Prepared by" table (and from undertaking.tex);
%%   - to show the co-supervisor or committee members 2 and 3, remove
%%     the leading % from those lines.
%% =====================================================================

\begin{titlepage}
\thispagestyle{empty}
\begin{center}

\vspace{4mm}
\includegraphics[width=0.5\textwidth]{CCSIT-Logo.png}

\vspace{6mm}

{\Large \textbf{\ksa}}\\[2mm]
{\Large \textbf{\moe}}\\[2mm]
{\Large \textbf{\kfu}}\\[2mm]
{\Large \textbf{\ccsitCollege}}\\[10mm]
{\large \textbf{\ccsitFull}}\\[10mm]
\vfill

{\LARGE \textbf{\gpTitle}\par}
\vspace{6mm}

{\large\itshape
\setstretch{1.05}
A project submitted\par
in partial fulfillment of the requirements for the degree of\par
\gpDegree\par
}

\vfill

\textbf{Prepared by}\\[2mm]
\begin{tabular}{rl}
\gpStudentA & (\gpStudentAID) \\
\gpStudentB & (\gpStudentBID) \\
\gpStudentC & (\gpStudentCID) \\
\end{tabular}

\vfill

\textbf{Supervised by}\\[2mm]
\gpSupervisor\\

\vspace{6mm}

\textbf{Committee Members}\\[2mm]
\gpCommitteeA\\

\vfill

{\gpSubmissionMonth}

\end{center}
\end{titlepage}

%% ----------------------------- Acknowledgment Page -----------------------------
\cleardoublepage
\thispagestyle{empty}
\begin{center}
{\Large \textbf{Acknowledgment}}
\end{center}

\vspace{6mm}
\gpAcknowledgmentText

%% ----------------------------- Undertaking -----------------------------
\cleardoublepage
\thispagestyle{empty}
\begin{center}
{\Large \textbf{Undertaking}}
\end{center}

\vspace{6mm}
\gpUndertakingText

\vfill
\noindent \textbf{Prepared by:}\\[3mm]
\begin{tabular}{@{}p{0.62\textwidth}p{0.34\textwidth}@{}}
\gpStudentA\ (\gpStudentAID) & Signature:\ \rule{45mm}{0.4pt}\\[4mm]
\gpStudentB\ (\gpStudentBID) & Signature:\ \rule{45mm}{0.4pt}\\[4mm]
\gpStudentC\ (\gpStudentCID) & Signature:\ \rule{45mm}{0.4pt}\\[4mm]
\end{tabular}

\vspace{6mm}
\noindent Date:\ \rule{55mm}{0.4pt}

%% ----------------------------- Front Matter -----------------------------
\cleardoublepage
\pagenumbering{roman}
\setcounter{page}{1}

\section*{Abstract}
\addcontentsline{toc}{section}{Abstract}

Scientific research documents convey important data both in their written form and through various graphical features, such as intricate figures, tables, and charts. Yet most conventional Retrieval Augmented Generation (RAG) systems focus mainly on text, leading to a substantial loss of the key information contained in these visual elements. Moreover, existing Artificial Intelligence (AI)-based document processing systems place a great deal of reliance on cloud-based Application Programming Interface (API), thereby creating serious risks in the areas of data privacy, protection of intellectual property, and the general reliability of the system. To overcome these important limitations, this project presents a thorough and privacy-preserving solution that uses AI to create a local multimodal repository of research documents. The system in question is intended to function completely offline, thus keeping sensitive academic and proprietary data secure while simultaneously carrying out a comprehensive analysis of both the text and the visual parts of scientific literature. In technical terms, the system functions by using deep layout-aware parsing and carries out separate extraction procedures for the textual and visual elements to maintain the original document structure. For the purpose of information retrieval, the architecture makes use of an advanced multimodal method that combines hybrid text search-this involving both keyword matching and dense semantic analysis-with independent visual search operations, the two different streams being smoothly integrated by means of the Reciprocal Rank Fusion (RRF) algorithm. To produce responses, the system uses a Vision‑Language Model (VLM) that is hosted locally and fine‑tuned on scientific multimodal datasets. This approach removes the need for any cloud services completely. The project contains three components. First, a functional web‑based repository makes document management easy. Second a powerful local AI processing pipeline guarantees offline operation. Third a thorough performance evaluation checks both retrieval and question‑answering accuracy. Finally, the project delivers an intelligent environment that allows researchers and educational institutions to query and understand complex scientific literature efficiently while keeping full control of their sensitive data.


\vspace{4mm}
\noindent\textbf{Keywords:} Research documents; Vision Language Model (VLM); Retrieval Augmented Generation (RAG); Local model.

\cleardoublepage
\tableofcontents

\cleardoublepage
\listoftables
\addcontentsline{toc}{section}{List of Tables}




\cleardoublepage
\listoffigures
\addcontentsline{toc}{section}{List of Figures}

%% ----------------------------- Acronyms -----------------------------
\cleardoublepage
\section*{List of Acronyms}
\addcontentsline{toc}{section}{List of Acronyms}

\begin{acronym}[123456]  
  \acro{API}{Application Programming Interface} 
  \acro{AI}{Artificial Intelligence}
  \acro{ASR}{Automatic Speech Recognition}
  \acro{BM25}{Best Matching 25}
  \acro{CCA}{complex concentrated alloys}
  \acro{FAISS}{Facebook AI Similarity Search}
  \acro{FTE}{full-time equivalent}
  \acro{GPU}{Graphics Processing Unit}
  \acro{HEA}{high entropy alloys}
  \acro{HITL}{Human-in-the-Loop}
  \acro{IR}{Information Retrieval} 
  \acro{LLM}{Large Language Model}
  \acro{LVLMs}{large vision language models}
  \acro{MAE}{mean absolute error}
  \acro{MDA}{Materials Database Agent}
  \acro{MIP}{Multimodal Ingestion Pipeline}
  \acro{MQA}{multimodal questions and answer}
  \acro{nDCG@5}{normalized discounted cumulative gain at 5}
  \acro{OCR}{Optical Character Recognition}
  \acro{OLAP}{Online Analytical Processing}
  \acro{PFTFI}{Fine Tuning with Feedback Inheritance}
  \acro{QA}{question-answering}
  \acro{RAG}{Retrieval Augmented Generation}
  \acro{RCG}{representation compression via generation}
  \acro{RCR}{representation compression via retrieval}
  \acro{RRF}{Reciprocal Rank Fusion} 
  \acro{TQA}{text-only questions and answer} 
  \acro{VLM}{Vision Language Model}
  \acro{VQA}{visual questions and answer}
  % \acro{VQA}{vision-only QA}
\end{acronym}

%% ----------------------------- Main Matter -----------------------------
\cleardoublepage
\pagenumbering{arabic}
\setcounter{page}{1}

\section{Introduction}
Research documents contain valuable information in both textual and visual forms, and nowadays, the use of \ac{RAG} systems is increasing rapidly. However, conventional \ac{RAG} systems primarily focus on understanding and retrieving textual information \cite{Tanaka2025VDocRAGRG}, resulting in loss of significant knowledge displayed through visual content, such as figures, diagrams, tables, and charts, in the retrieval process. This limitation affects the completeness of retrieved information and the quality of the generated results.\\

To address this limitation, our project aims to apply local multimodal retrieval techniques that ensure visual content is considered in information retrieval, which will support a wider range of research document questions and improve the quality of generated results. To provide a clear overview of the project, we will start by discussing fundamental components relevant to the project followed by its core objectives. Then, we will explain the importance of the project, which will provide the context for the specific problem we want to address. The last aspect to be discussed in this section is the project boundaries and expected outcomes.

\subsection{Background}
Processing multimodal documents requires specialized pipeline to enhance the efficiency of the use of recent technologies in this field. The pipeline mainly relies on four interconnected techniques that are fundamentally relevant to context search and answer generation, which are \ac{IR}, \acp{LLM}, \acp{VLM}, and Retrieval Augmented Generation (\ac{RAG}). 

\subsubsection{Information Retrieval}
\ac{IR} involves arranging, searching, and retrieving useful information from a pool of data in response to a query made by a user. Traditional \ac{IR} methods rely on calculating the degree of relevance of documents through techniques like term frequency, inverse document frequency, and vector space model \cite{Bhati2026IntelligentIR}. The more advanced intelligent retrieval systems, however, leverage deep learning and semantic modeling to capture context and the connection between the queries and documents for effective \ac{IR}.\\

In the AI-Based Research Documents Repository Using \ac{LLM}, the process of \ac{IR} can be applied to look through the collection of stored research documents and find relevant information for the user's query. Rather than forcing the user to go through numerous research papers themselves, the system can represent both the query and the documents, calculate their similarity and relevance, and obtain the most relevant research information. This information will help the \ac{LLM} generate responses, which will derive from the knowledge contained in the research repository. This way, \ac{IR} becomes the link between the user's query, the stored research documents, and the \ac{LLM}.

\subsubsection{Large Language Models}
\acp{LLM} are sophisticated deep learning systems that are intended to understand, produce, and manipulate human language. By using massive amounts of textual data for training, these models apply transformer architectures to pick up the complex semantic and contextual relationships present in the data \cite{Zhao2023ASO}. \acp{LLM} are usually divided into closed-source ones, for example GPT-4, and open-source ones, such as Llama and Qwen. In academic and secure settings, open-source models are becoming more popular because of their cost efficiency, better data privacy, and the fact that they can run on local infrastructure. Although they have strong language capabilities, \acp{LLM} are limited by their fixed training data, potentially causing knowledge gaps or hallucinations \cite{Ji2022SurveyOH}. To overcome these drawbacks and increase factual accuracy, current \ac{AI} systems often combine \acp{LLM} with external knowledge retrieval systems, such as \ac{RAG}, so that the models can create responses that are properly grounded in the context.

\subsubsection{Vision Language Models}
\acp{VLM} are multimodal AI models that combine visual and linguistic capabilities through the integration of data from images and videos with natural language. Such models are trained using methods like masked modeling, which helps predict missing parts of an image or text, as well as image-text matching, which enables the identification of an association between a given description and an image. As shown in Figure~\ref{fig:VLM-architecture}, an important step in \acp{VLM} is embedding, where textual and visual inputs are processed by text and vision encoders and transformed into text embeddings and image embeddings. These embeddings provide numerical representations of the input data, allowing the model to capture and relate information from both visual and textual modalities. These techniques have resulted in the emergence of state-of-the-art models such as Gemini, OpenFlamingo, and Goldfish \cite{Bordes2024AnIT}, as well as model architectures like unified multimodal models (Chameleon), two-tower models (CLIP), and encoder-decoder models (BLIP-2) \cite{Ray2024ERVQAAD}, as shown below in Figure~\ref{fig:VLM-architecture}. The figure below represents the \ac{VLM} architecture used to illustrate two commonly used \ac{VLM} techniques. In both methods, visual and text inputs are mapped into visual and text embeddings, respectively. This enables the \ac{VLM} to convert the information into numeric form.

\begin{figure}[H]
\centering
\includegraphics[width=1\textwidth]{VLM Architecture.png}
\caption[Architectural components of modern \acsp{VLM}]{Architectural components of modern \acp{VLM} \cite{Li2025ASO}.}
    \label{fig:VLM-architecture}
\end{figure}

\subsubsection{Retrieval Augmented Generation}
 \ac{RAG} is a framework that supports and enhances information retrieval and language generation. \ac{RAG} is used to address key issues faced by \acp{LLM}, such as hallucination \cite{Huang2024ASO}. Hallucination arises due to the \acp{LLM} reliance on the training data. The \ac{RAG} workflow addresses this issue in three main stages \cite{Huang2024ASO}. The first stage is indexing, which consists of data preparation to facilitate the search in the next stage. The second stage is retrieval, which is the core stage in most systems. Retrieval is done using either traditional search methods or current search methods. Traditional search methods focus on term frequency; however, current search methods rely on semantic search, which may retrieve more relevant information. The last stage is generation, which combines the query with retrieved context. The common method used involves simple concatenation of the query question and retrieved information; the result is then supplied to the \ac{LLM} to generate the answer.
 
\subsection{objectives}
Our project considers several aspects relevant to multimodal retrieval. To handle these aspects, we have to use combination of several steps and tools. Together, they build the foundation of our project. The project aims to achieve mainly four objectives in which they specify core aspects of the project that are used together to address the project problem. The objectives are explained as follows:\\

O1: To develop a functional web-based \ac{AI} research documents repository system that enables users to successfully upload and manage documents, with a success rate of at least 90\%, evaluated by Week 7 of the second semester.\\

O2: To implement a local multimodal \ac{AI} pipeline that processes and retrieves textual and visual information from research papers, achieving retrieval accuracy of at least 60\%, evaluated by Week 10 of the second semester.\\

O3: To develop a fine-tuned local language model to answer users’ questions based on textual and visual information, achieving question answering accuracy of at least 65\%, evaluated by Week 12 of the second semester.\\

O4: To integrate all system components into a unified system that enables multimodal question answering for research documents, generating answers with an accuracy of 65\%, evaluated by Week 13 of the second semester.

\subsection{Motivation}
The problem that both students and researchers face is the massive amount of scientific research work, making it very hard to find and understand the information because of its scattered presence in many documents. Research papers may contain information that goes beyond just text; they might include graphs, tables, images, and more. Traditional databases and search engines cannot make sense of such multimodal information, and thus the user has to go through many documents. There is a definite need for an intelligent research document repository that understands text as well as visuals and helps in \ac{IR}.\\

In addition, current document processing systems utilizing \ac{AI} rely heavily on cloud services and external \ac{API}. This raises concerns regarding document privacy, internet reliability, and external data processing; consequently, the use of local servers is crucial for ensuring security. Processing language models and \acp{VLM} locally enable document processing, retrieval, and question-answering to occur entirely within the system itself. Accordingly, multimodal document understanding technologies can be integrated with local \ac{RAG} to create a system that is both more secure and better equipped to comprehend context.

\subsection{Problem Statement}
Even with the latest developments in document question answering and multimodal \ac{RAG}, there are a few critical challenges posed by current techniques in dealing with complicated visual research documents. The first challenge is the heavy dependence of visual retrieval on \ac{OCR}-aligned data, which makes the extraction and retrieval of non-textual visual data such as diagrams, raw charts, and figures very challenging. Another challenge posed by alternative deep research techniques is the use of multiple retrieval and refinement loops per query. Furthermore, standard fixed-size data chunking regularly serves related ideas across arbitrary boundaries, while current hybrid retrieval methods combine keyword and semantic search solely on textual chunks without integrating independently retrieved visual evidence.\\

Consequently, there is a distinct gap in the availability of an effective local, privacy-preserving system capable of preserving document layout structure, independently processing textual and visual modalities, and fusing evidence from both streams during retrieval. This project addresses this issue by proposing a pipeline for an offline multimodal \ac{RAG} system. The proposed architecture performs deep layout-aware document structure analysis, employs hybrid text retrieval (combining the \ac{BM25} algorithm with dense semantic search) alongside independent visual retrieval, and unifies multimodal results at the retrieval stage using the \ac{RRF} algorithm all prior to response generation by a locally hosted \ac{VLM}.

\subsection{Scope and Deliverables}
This section is devoted to defining the functional and technical boundaries within which the project operates by accurately defining the scope of work. This section reviews the compatibility of the project with the strategic identity of King Faisal University and its societal impact and concludes with a statement of the concrete outputs expected at the end of the project period to prove the effectiveness of the achieved system.

\subsubsection{Definition of the Project's Scope}
Scope refers to the functional and technical limitations of the system, as it is constrained by its capability to process only research and technical documents in PDF format to extract and analyze textual and visual information, such as tables, figures, and graphs. It will assist in the accurate search, retrieval, and summary of information using a multimodal retrieval enhanced generation approach. The limitation of the system's processing ability includes local processing for confidentiality. The system is created specifically for the use of researchers, students, and educational institutions for the purpose of analyzing their research documents privately and confidentially.

\subsubsection{Alignment with KFU Identity and Society Impact}
The proposed AI-Based Research Documents Repository Using \acp{VLM} aligns with King Faisal University's trend toward digital transformation, scientific research, innovation, and the effective use of modern technologies. KFU emphasizes the use of technology to support teaching, learning, and research services, as well as improving access to knowledge resources and research information.\\

The proposed approach to achieve these goals will be through offering an intelligent research repository that uses multimodal document processing, hybrid retrieval, and local \ac{AI} models to allow for efficient and interactive access to research knowledge for students and researchers. This approach will enable the system to obtain data from both visual and textual content in the research documents and produce grounded answers.\\

In terms of societally focused applications, the implementation of the project would allow easier access to scientific knowledge, as well as reduced time and effort involved in researching vast sets of literature materials. Moreover, the technology implemented by the project can make users more independent regarding their control over research data by decreasing reliance on \ac{AI} services located on cloud networks. Thus, the contribution of the project to form a knowledge-based, digitally advanced society is clear. In addition, the project is directly associated with the development activities of KFU.

\subsubsection{Expected Project Deliverables}
The expected outputs include the delivery of a functional web-based research repository that allows uploading and managing research papers, provides multimodal search and retrieval capabilities, answers queries in natural language, summarizes and analyzes documents, and integrates the processing of textual and visual content. The project includes a fully operational \ac{AI} processing pipeline in a local environment, ensuring local processing of documents, multimodal retrieval, and local generation without relying on cloud services. The output also includes a specially trained and adapted local language model to answer the questions of research documents with high accuracy. A comprehensive evaluation results package will be presented that measures the efficiency of the system by analyzing the retrieval performance and the accuracy of answering questions to demonstrate the effectiveness of the proposed methodology and its suitability for the research environment.

\cleardoublepage
\section{Literature Review}
In this section, we will analyze some previous studies that have been done in the field of multimodal document research. We will discuss important aspects of each study including the methodology, datasets if any, and final results. We have collected the most relevant researches to our topic to identify their research gaps that will help us to build our project.\\

Tanaka et al. \cite{Tanaka2025VDocRAGRG} (2025) developed a new \ac{RAG} framework, VDocRAG, which answers questions over visually rich documents. Unlike text-based \ac{RAG}, VDocRAG treats the whole document as an image instead of converting it completely to text. The framework consists of two main components. First, VDocRetriever is responsible for retrieving relevant document images. Second, VDocGenerator is responsible for generating the answers based on the retrieved document images. To improve visual document retrieval process, Tanaka et al. \cite{Tanaka2025VDocRAGRG} adopted pre-trained \ac{LVLMs} and introduced two self-supervised pre-training tasks. First, \ac{RCR} associates document images with their \ac{OCR} text. Second, \ac{RCG} trains the model on reconstructing the \ac{OCR} tokens from the visual representations.\\

Tanaka et al. \cite{Tanaka2025VDocRAGRG} introduced OpenDocVQA, which is a collection of refined \ac{VQA} datasets containing visually rich documents, including tables, slides, charts, webpages, infographics. The evaluation of models is done with both single-pool and all-pool settings. The single-pool setting specifies that documents are retrieved from each dataset separately, while the all-pool setting requires retrieving documents from all datasets. Retrieval performance is evaluated using \ac{nDCG@5}, which measures the quality of the top 5 retrieved results. VDocRetriever achieved \ac{nDCG@5} scores of up to 86.0 in the single-pool setting and 76.4 in the all-pool setting. VDocRAG outperformed text-based \ac{RAG} across all datasets used for evaluation. For example, text-based \ac{RAG} achieved 28.0\% relaxed accuracy on ChartQA in single-pool settings, while VDocRAG achieved 52.0\%. The results demonstrate that representing documents directly as images enhances the performance of question answering over visually rich documents.\\

Dong et al. \cite{Dong2026DocResearcherAU} (2026) developed a new system called Doc-Researcher for answering complex questions over collection of documents containing text with figures, tables, charts, and other elements. The methodology is divided into two stages. The first stage is the offline stage or deep document parsing and indexing. It applies layout-aware parsing to each document for extracting document’s elements, such as text, tables, charts, figures, and equations. Each element is labeled with its layout type with a bounding box as shown in Figure~\ref{fig:layout-aware_chunking}. Then, a one-time textual conversion is performed on visual content. Qwen2.5-VL is a large vision-language model, it was used for generating overall summaries and detailed descriptions of figures and tables. Parsing is done to avoid repeatedly interpreting the same content for every question.

\begin{figure}[H]
\centering
\includegraphics[width=1\textwidth]{layout-aware chunking.png}
\caption[Overview of layout-aware chunking process]{Overview of layout-aware chunking process \cite{Dong2026DocResearcherAU}.}
\label{fig:layout-aware_chunking}
\end{figure}

The second stage of the methodology is the online stage or Doc-Researcher framework for multimodal deep research. It organizes the research process into four roles. First, the planner divides the query into sub-queries Second, the searcher retrieves the relevant documents identified by planner. Third, the refiner removes duplicates and assesses whether the collected evidence is sufficient and initiates additional retrievals when necessary. Fourth, the reporter generates the answer to the question by combining the collected evidence alongside the citation metadata. Doc-Research has achieved an overall answer accuracy of 50.6\%, While the hybrid retrieval, which combines text-only retrieval and vision-only retrieval, achieved 42.4 - 50.6\% accuracy, which outperforms other retrieval techniques by 3 - 5\%.\\

Gosmar and Zenezini \cite{Gosmar2026MADPAM} (2026) introduced MADP, a multiagent architecture designed to automate enterprise document processing by integrating deep learning for classification and analysis with \ac{LLM} for data mining, while ensuring high accuracy through selective human verification. The methodology comprises five sequential modules: a Classifier Agent utilizing a ResNet-18 CNN to identify document types and suppliers; a Splitter Agent to separate logical units in multi-page documents; a Parser Agent using the Docling library to convert formats to Markdown, thereby reducing noise and tokens; an Extraction Agent employing \ac{LLM} with specific prompts to deduce JSON information; and a Validator Agent applying a novel \ac{PFTFI} mechanism for post-processing. This pipeline is supplemented by a \ac{HITL} interface, where human moderators correct extracted metadata.\\

Evaluated using document-level accuracy, field-level F1 scores, and processing time, the system was tested on a dataset of 955 multilingual (Italian, English, German, French, and Arabic) documents from over 50 heterogeneous suppliers, alongside a 100-document stratified subset for ablation studies and a projected scenario of 100,000 annual invoices. In production, the system achieved a 97.0\% full automation rate, requiring non-\ac{AI} intervention for only 3\% of documents. The ablation study showed that MADP with \ac{HITL} supervision reached a 98.5\% document-level accuracy, with the Parser Agent alone improving accuracy by 17.5 percentage points. Model evaluations revealed that Mistral-Small-3.2 achieved the highest F1 score (92.9\%), while DeepSeek-OCR-Regolo offered the fastest inference time. Scaling to 100,000 invoices annually reduces \ac{FTE} labor requirements by nearly 70\%, and the hybrid \ac{AI}+\ac{HITL} approach significantly boosts sustainability by cutting CO2 emissions and energy consumption by 69\%, and water usage by 63\% compared to traditional manual processing.\\

Chandrasekhar et al. \cite{Chandrasekhar2026MaterialDA} (2026) introduced the \ac{MDA}, a multi-agent system architecture designed to automate the extraction of structured data from materials science literature by converting research PDFs into unified tabular databases using multimodal \ac{LLM} and parallel processing. The \ac{MDA} methodology comprises four sequential stages: an input and extraction stage that processes PDFs using the MCP model and the marker-pdf tool to generate Markdown files for text and tables, and JPEG images for graphs; a data decomposition stage that organizes outputs into local subfolders per paper to preserve multimodal context; a parallel processing stage where a main agent spawns independent "doc-writer" sub-agents to synchronously read subfolders and extract material parameters into structured JSON files; and a compilation stage where a "csv-writer" agent unifies the JSON files into a single CSV database. The system is prompt-driven and incorporates \ac{HITL} supervision for guidance and verification.\\

The system was evaluated using two experimental MeltpoolNet manufacturing kits containing 789 rows of process parameters and material properties, alongside a dataset of \ac{HEA} and \ac{CCA} comprising 366 rows of mechanical properties embedded within graphs, diagrams, and stress-strain curves rather than regular text. Results demonstrated that \ac{MDA} is a highly scalable solution with performance varying by data type. For the text and table-rich MeltpoolNet dataset, the GLM 5V Turbo model achieved the highest F1 scores, while Qwen-3.5 recorded the lowest \ac{MAE} for melting temperature. Conversely, for the graph-heavy \ac{HEA}/\ac{CCA} dataset, Claude Opus 4.6 outperformed other models by accurately reproducing mechanical properties with minimal errors, whereas GPT Series models exhibited larger errors. Ultimately, the GLM 5V Turbo model was identified as the most balanced base model for handling diverse data formats.\\

A. R and R. Dk \cite{R2026DOCGA} (2026) proposed DOC GPT, an offline document-based chatbot that answers user questions according to uploaded PDF documents, without using public \ac{API} services. The methodology is based on \ac{RAG} and consists of the following operation sequence. Initially, a PDF file is uploaded by the user, and information is extracted from the PDF file using \ac{OLAP} and \ac{OCR}, then the information is divided into chunks while retaining context. Third, these chunks are encoded into semantic vectors via Sentence Transformers and saved in the vector DB (e.g., \ac{FAISS}, Chroma). Fourth, in the case of user's question the system performs a semantic similarity search and reranks of the most relevant chunks. Fifth, this context is passed to a local \ac{LLM} via Ollama (Llama, Mistral, etc.) to produce grounded answers. Sixth, the generated answer is returned to the frontend and shown to the user. The use of local machines allows the system to maintain document privacy and operate offline, while also enabling web search when necessity.\\

To perform the system evaluation, two major datasets were utilized: DocVQA and the ArXiv Papers Dataset. DocVQA contains a collection of thousands of document images such as scans of reports, letters, and multi-column forms accompanied by the corresponding natural-language questions and ground-truth answers. It was used to evaluate the ability of the system's work with complex layouts, scans, visual elements, and \ac{OCR} text. The second dataset, the ArXiv Papers Dataset, comprises a large set of scientific papers related to various research areas. This dataset helped to test the ability of the system's work with long documents, complex technical language, text chunks, and semantic search. Based on the reported results, the system demonstrated very good response accuracy and did not hallucinate, as it was strongly grounded in the locally retrieved document contents. Thanks to the combination of \ac{FAISS} vector search and reranking, it was possible to provide high-precision document retrieval, while \ac{OCR} worked successfully with complex, scanned, and multicolored document layouts. Ollama facilitated the entire process of work locally and offline.\\

Ahamed et al. \cite{AlthafAhamed2026AMA} (2026) proposed "A Multimodal and Multilingual Offline Retrieval-Augmented Generation System Using Local \ac{LLM}," an approach that is offline and privacy-preserving for deploying \ac{RAG}-based systems without using any \ac{API} or internet connection. This approach works on the \ac{MIP} framework , which accepts various inputs like PDFs, scanned documents, images, and audio recordings. Digital PDF inputs are converted to text using PyMuPDF. Scanned documents and images are processed using Tesseract \ac{OCR}. Audio recordings are transcribed using \ac{ASR}. The text obtained from different types of inputs is further segmented into small fragments, and dense embeddings are created. A Hybrid Search method performs both keyword-based retrieval and dense retrieval of the most semantically similar content using a combination of \ac{BM25} and ChromaDB. The retrieved results are merged, ranked using \ac{RRF}, and the obtained context is passed as input to the local Mistral-7B through Ollama to get citation-backed responses. As illustrated in Figure~\ref{fig:Fusion-Seek}, Fusion-Seek implements a method based on the functioning of four different layers, namely multimodal input, text extraction, chunking, embedding creation, hybrid retrieval, and finally, response generation using a local \ac{LLM}.

\begin{figure}[H]
\centering
\includegraphics[width=0.5\textwidth]{Fusion-Seek.png}
\caption[Fusion-Seek: Four-Layer End-to-End System Architecture]{Fusion-Seek: Four-Layer End-to-End System Architecture \cite{AlthafAhamed2026AMA}.}
\label{fig:Fusion-Seek}
\end{figure}

This work was implemented and evaluated based on technologies such as PyMuPDF, Tesseract \ac{OCR}, \ac{ASR}, \ac{BM25}, ChromaDB, Ollama, and Mistral-7B. The proposed approach can process various types of documents such as technical manuals, policy documents, scientific papers, scans and images, and meeting transcripts. These experiments showed that the combination of \ac{BM25} keyword search, dense vectors, and \ac{RRF} can increase the quality of contextual search and decrease hallucinations in contrast to vector-based-only search. Using Ollama and Mistral-7B proved that it is possible to generate citation-grounded responses completely offline.\\

Chen et al. \cite{Chen2026SciMDRAS} (2026) introduced SciMDR, which is a \ac{QA} dataset designed to train models for understanding multimodal content within research documents, such as text, figures, and tables. It consists of 304,461 \ac{QA} pairs derived from approximately 20,000 research papers; the \ac{QA} pairs are divided into three types: \ac{TQA} pairs, vision-only QA pairs (\ac{VQA}), and \ac{MQA} pairs that require information from both text and visuals. The dataset was constructed in two stages. First, research papers were parsed with MinerU2.5 \ac{OCR} to extract text, captions, figures, and tables. The text information was then fed to GPT-5.1 to extract discrete claims from the text. These claims were used to guide the generation of questions. By supplying the claims, the model generates questions whose answers correspond to the provided claims. Then, claims with their associated visuals are used to generate \ac{MQA} pairs, while text-only claims are used to generate \ac{TQA} pairs. \ac{VQA} pairs are generated separately by providing visual information to the model.\\

Second, the generated \ac{QA} pairs are used with their original document context to support model learning to locate relevant evidence and reason over it when answering questions. Chen et al. \cite{Chen2026SciMDRAS} also introduced the SciMDR-Eval benchmark, which was manually created from 300 research papers. It consists of 907 \ac{QA} pairs, where the questions require models to search for and reason over information available in the full document, such as text, figures, tables, or a combination of them. The benchmark evaluates the model’s ability to locate relevant evidence, reason over it, and generate correct answers. The performance of models fine-tuned on SciMDR was evaluated on various benchmarks, such as ChartQA, CharXiv, SPIQA, and SciMDR-Eval. The results showed that fine-tuning the Qwen2.5-VL-7B model on the SciMDR training dataset increased its SciMDR-Eval accuracy from 19.8\% to 49.1\%, an improvement of 29.3 percentage points.\\

The summary of literature review is represented in Table~\ref{tab:Literature}, which outlines the main aspects of each research paper, including the methods used, datasets applied, and achieved performance. Following the summary table, next chapter will discuss the proposed methodology and demonstrate how it can overcome previous studies challenges.

\begin{table}[h]
    \centering
    \begin{tabular}{|>{\raggedright\arraybackslash}p{0.14\textwidth}|>{\raggedright\arraybackslash}p{0.23\textwidth}|>{\raggedright\arraybackslash}p{0.23\textwidth}|>{\raggedright\arraybackslash}p{0.26\textwidth}|}
        \hline
        \textbf{Authors} & \textbf{Method} & \textbf{Dataset} & \textbf{Performance} \\
        \hline
        \textbf{Tanaka et al.} &
        VDocRetriever, VDocGenerator, \ac{RCR}, \ac{RCG} &
        OpenDocVQA &
        Up to 86.0 \ac{nDCG@5} for retrieval, 52.0\% relaxed accuracy on ChartQA for generation \\
        \hline
        \textbf{Dong et al.} &
        Deep Document Parsing, Multimodal Deep Research, Hybrid Retrieval &
        M4DocBench for evaluation &
        50.6\% answer accuracy on M4DocBench and 3.4x higher than state-of-the-art baselines \\
        \hline
        \textbf{Gosmar and Zenezini} &
        MADP, \ac{HITL}, \ac{PFTFI} &
        955 production documents; 100-document evaluation subset &
        98.5\% accuracy with \ac{HITL}, compared with 85\% for pure \ac{AI} automation \\
        \hline
        \textbf{Chandrasekhar et al.} &
        \ac{MDA}, MarkerPDF &
        MeltpoolNet and \ac{HEA}/\ac{CCA} databases for evaluation &
        Up to 75.00\% F1 for figure-derived depth-of-MeltPool extraction \\
        \hline
        \textbf{R. and Dk.} &
        \ac{RAG}, \ac{OCR} &
        DocVQA and arXiv Papers Dataset for evaluation &
        Improved semantic retrieval accuracy; no numerical accuracy reported \\
        \hline
        \textbf{Ahamed et al.} &
        \ac{MIP}, \ac{OCR}, \ac{ASR}, Hybrid Search (\ac{BM25} and semantic search), Ollama &
        Technical manuals, legal/policy documents, research papers, images, scanned documents, audio &
        87.4\% Recall@10, 84.9\% Precision@10, 0.82 MRR with hybrid \ac{RRF} \\
        \hline
        \textbf{Z. Chen et al.} &
        Scientific multimodal dataset construction &
        SciMDR for training and SciMDR-Eval for evaluation &
        49.1\% accuracy on SciMDR-Eval when Qwen2.5-VL-7B is fine-tuned on SciMDR, which is 29.3\% higher \\
        \hline
        \textbf{Proposed Strategy} &
        A combination of Deep Document Parsing, Hybrid Search (\ac{BM25} and semantic search), \ac{RAG}, and Hybrid Retrieval &
        SciMDR~\cite{Chen2026SciMDRAS} &
        -- \\
        \hline
    \end{tabular}
    \caption{Summary of literature review}
    \label{tab:Literature}
\end{table}

\cleardoublepage
\section{Methodology}
In this section, we will discuss our methodology by establishing a foundation for our selection of the methodology. We will begin by addressing the key challenges associated with methodologies used in previous studies. Then, we will propose a strategy for each challenge to overcome it. Next, we will justify our selection clearly to highlight the main aspects of our methodology, and we will explain how it may overcome the challenges previously explained. Finally, we will discuss the tools to be used to perform our methodology in the most efficient way.

\subsection{Challenges and Proposed Strategies}
Previous studies have implemented various techniques in evidence retrieval and search. However, each technique had its gaps, we went through them and selected the most relevant challenges to our project. This section will discuss each challenge and a proposed strategy to overcome it.

\subsubsection{Challenge 1: Limited Retrieval of Textless Visual Information}
VDocRAG \cite{Tanaka2025VDocRAGRG} uses supervised pre-training to align visual content and \ac{OCR}-extracted text data for document retrieval. This approach allows the model to effectively retrieve information from visually rich documents when the visual content contains text. However, the reliance on the alignment between visual content and \ac{OCR}-extracted text data limits the model’s ability to retrieve information from visuals that do not contain text. Moreover, the model is designed primarily for visually rich documents, which can degrade the performance of VDocRAG when dealing with text-heavy documents.\\

\textbf{Proposed strategy:} To fix this issue, the system processes content independently without relying on text extracted through \ac{OCR}. Instead of using text to match visuals, it employs a special visual model called ColQwen2. This model examines image pixels directly and creates visual embeddings. That way, even if an image lacks text, the system can still find visual information based on what the image shows. 

\subsubsection{Challenge 2: High Computational Cost}
Doc-Researcher \cite{Dong2026DocResearcherAU} uses a deep research approach that increases computational costs. The deep research approach requires performing evidence retrieval and refinement iteratively until a predefined condition is met. During this process the system iteratively retrieves and evaluates evidence to maximize the sufficiency of the collected evidence for answering the user question. Therefore, multiple evidence retrieval and refinement iterations may be performed for each query. The iterative process can increase the overall computational costs and processing time, especially when dealing with large collection of complex multimodal documents.\\ 

\textbf{Proposed strategy:} To reduce computing demands, the system uses a simpler, non-iterative workflow. It processes the document in one pass-parsing, embedding, retrieving, and generating-without looping back multiple times for each query. By skipping repeated retrieval and refinement steps, the process becomes much faster and uses fewer resources while still providing good context for answering questions. 

\subsubsection{Challenge 3: Limited Multimodal Information Extraction}
The \ac{MDA} system \cite{Chandrasekhar2026MaterialDA} primarily processes the document into machine-readable content. The MDA system uses an extraction approach that is limited to textual information embedded in the document. Therefore, the system may not fully capture the scientific information represented in visual content, as its main process is the processing of captions and paragraphs relevant to images. As a result, important measurements and other scientific information contained within images may be lost or not fully captured, leading to incomplete representations of research documents, particularly when they contain essential scientific information embedded in images.\\

\textbf{Proposed strategy:} To avoid losing details, the system employs a deep multimodal parsing step with the MinerU library. Rather than just grabbing captions or nearby text, it extracts actual visual components such as graphs, charts, and figures. These are then analyzed using Qwen2.5-VL-7B-Instruct. This allows the system to grasp the meaning within images, leading to a complete and accurate representation of the document's content.

\subsubsection{Challenge 4: Limitations in Context Retrieval}
DOC GPT \cite{R2026DOCGA} uses static-size chunking when processing textual information. This type of chunking might potentially divide connected ideas into different chunks and could lead to limited context and missing complete information about the query. In addition, chunks may include two or more ideas that might be irrelevant to the intended query.\\

\textbf{Proposed strategy:} To keep ideas connected and prevent context loss, the system replaces fixed-size text chunks with layout-aware chunking. It reads the document's structure and splits the text into meaningful units, such as paragraphs, sections, or blocks, based on how the document is organized. This ensures that each chunk retains coherent thoughts and does not mix unrelated parts.

\subsubsection{Challenge 5: Limitations in Cross-Modal Retrieval}
The Multimodal and Multilingual Offline \ac{RAG} System \cite{AlthafAhamed2026AMA} applies a reranking technique called RRF that combines the keyword and semantic searches for only document text chunks. However, the system cannot combine and rerank across multimodal evidence. As a result, the system may struggle to combine relevant text and images with the user’s query. \\

\textbf{Proposed strategy:} To improve cross-modal retrieval, the system adds a multimodal hybrid retrieval stage. After searching for text and visuals, it generates two ranked lists-one for text chunks and another for images. Then, it applies the \ac{RRF} algorithm to combine both lists. This fusion occurs at retrieval time, merging text and image results into a single list. This helps the system gather multimodal evidence to answer the user's questions more effectively.


\subsection{Justifications of the Selected Approach}
Based on the methodologies discussed in previous studies, we are proposing two pipelines to address challenges faced by previous research. Both pipelines support multimodal \ac{RAG} development, which enhances local models’ ability to retrieve relevant information from both text and visual content in research documents. Each proposed pipeline uses a slightly different approach to process and retrieve relevant information. The pipelines are shown in Figure~\ref{fig:RAG-pipelines-representation} and explained as follows:\\ 

\begin{figure}[H]
\centering
\includegraphics[width=1\textwidth]{RAG pipelines representation.png}
\caption[Alternative and Proposed \acs{RAG} pipelines representation]{Alternative and Proposed \acs{RAG} pipelines representation.}
\label{fig:RAG-pipelines-representation}
\end{figure}

\subsubsection{Alternative Multimodal RAG Pipeline}
The pipeline consists of four stages. First, multimodal document parsing is the foundational stage in the pipeline. The document is preprocessed by extracting textual information, which includes \ac{OCR} extraction in the case of scanned documents. Second, the multimodal document representation stage extracts visual information by using a \ac{VLM} that generates descriptions for figures and tables to create a multimodal document representation. Moreover, the extracted information is divided into smaller chunks, which are converted into semantic embeddings and indexed locally for retrieval.\\

Third, the hybrid retrieval stage is mainly a strategy for the information retrieval process. The hybrid retrieval strategy is a combination of \ac{BM25} keyword retrieval and dense semantic retrieval to enhance the accuracy of retrieving information from the research documents. Using \ac{BM25}, the relevant information is retrieved through keyword matching between the user’s query and the text available in storage. On the other hand, dense semantic retrieval processes the query and the chunks of documents into vectors and analyzes their similarity in meaning. The results from \ac{BM25} and dense semantic retrieval are then merged using the \ac{RRF} method. The most highly ranked chunks are chosen to be delivered to the local language model as contextual information.\\

Finally, at the local generation stage, all information and evidence relevant to the question have been retrieved, ranked, and prepared for the final generation step. The local generation model can be executed locally using Ollama. Moreover, the model is provided with a prompt containing the user’s question and relevant \ac{RAG} context to be used in answer generation. The \ac{RAG} context consists of the top-ranked results of the hybrid retrieval output, including text chunks and visual descriptions. Therefore, the user’s question and the retrieved information shape the prompt, allowing the local \ac{LLM} to use the retrieved information as evidence when generating the answer.

\subsubsection{Proposed Multimodal RAG Pipeline}
The pipeline consists of four stages. First, deep multimodal document parsing is the foundational stage that helps to understand the document structure and visual semantics. We have adopted the Doc-Researcher’s deep multimodal document parsing \cite{Dong2026DocResearcherAU}, which specializes in multimodal document analysis, performing a deep analysis of multimodal documents and converting them into organized and searchable units while maintaining multimodal semantics. It divides the page into independent areas, identifying and extracting structural elements such as tables and figures, while accurately retaining their spatial coordinates and their contextual relationships within the page. Second, the multimodal document processing and indexing stage extracts visual information similar to the alternative multimodal \ac{RAG} pipeline by using a \ac{VLM} that generates descriptions for figures and tables. Then, the extracted textual and visual information is processed and indexed locally for retrieval.\\

Third, the hybrid multimodal retrieval stage is the strategy used for the information retrieval process. It begins with a hybrid text retrieval strategy, which is a combination of \ac{BM25} keyword retrieval and dense semantic retrieval to enhance the accuracy of retrieving information from the research documents. In addition, hybrid multimodal retrieval combines hybrid text retrieval results with visuals retrieved independently. The actual visuals are retrieved and ranked independently using a visual retrieval model, which allows the system to preserve visual information that may not be fully captured by textual descriptions. Consequently, we will have two independent ranked lists, one for text chunks and one for visuals. Then, retrieval-time multimodal fusion is applied to combine the rank positions produced by the two independent retrieval systems by using \ac{RRF}. In that process each visual and its associated description are treated as a single retrieval item, and the ranked list consists of both visual identifier and its associated description, allowing the system to retrieve the original visual at the generation stage. The resulting list represents the final hybrid multimodal retrieval output, which will be used to construct the final \ac{RAG} context.\\

Finally, at the local generation stage, all information and evidence relevant to the question have been retrieved, ranked, and prepared for the final generation step. The local generation model is fine-tuned with the SciMDR dataset \cite{Chen2026SciMDRAS}. After fine-tuning, the local generation model can be executed locally using Ollama. Moreover, the model is provided with a prompt containing the user’s question and relevant \ac{RAG} context to be used in answer generation. The \ac{RAG} context consists of the top-ranked results of the final hybrid multimodal retrieval output, including text chunks and figures with their associated descriptions. Therefore, the user’s question and the retrieved information shape the prompt, allowing the local \ac{VLM} to use the retrieved information with original figures as evidence when generating the answer.\\

The proposed multimodal \ac{RAG} pipeline is selected because it preserves the independence of textual and visual information during retrieval, and it only combines the results at the end of the retrieval process. Prepared context and original visual content are supplied to the local \ac{VLM} to generate the answers. On the other hand, the alternative multimodal \ac{RAG} pipeline simplifies the retrieval by relying on the conversion of visual content into descriptions before retrieval without independently retrieving the actual visuals. Moreover, only textual information and visuals descriptions are supplied to the local \ac{LLM}. As a result, the alternative multimodal \ac{RAG} pipeline may fail to retrieve visual content whose relevant information was not captured by the generated descriptions.

\subsection{Discussion of Tools and Techniques to be Used}
This project aims to integrate various technologies for document processing, visual understanding, search, and local language models to develop an \ac{AI}-based research documents database. The MinerU library \cite{Wang2024MinerUAO} will be used to perform layout-aware PDF parsing tasks, whereas PyMuPDF \cite{artifex2026pymupdf} will facilitate the handling of PDF and page-level tasks like extraction and rendering of pages. Tesseract \ac{OCR} \cite{Smith2007AnOO} will perform the task of extracting text from scanned pages or images. To process visual information (images, figures, etc.), Qwen2.5-VL-7B-Instruct \cite{Bai2025Qwen25VLTR} will be used to understand the information that cannot be comprehended using text extraction techniques alone. The use of visual information alongside textual content is also supported by multimodal document-processing approaches presented in \cite{Tanaka2025VDocRAGRG} and \cite{Chandrasekhar2026MaterialDA}.\\

To perform information retrieval, the semantic embedding of the meaning of the content in the documents will be created by BGE-M3 \cite{Chen2024M3EmbeddingMM} and stored/retrieved via ChromaDB \cite{huber2026chroma}. The \ac{BM25} algorithm \cite{Robertson2009ThePR} will also be used in conjunction with semantic search in the case of exact keyword matching. The ColQwen2 model \cite{Faysse2024ColPaliED} is used to retrieve visuals independently. The outputs of the text and visual searches will be merged and ranked using the \ac{RRF} algorithm \cite{Cormack2009ReciprocalRF}, which is an algorithm and not an independent software tool. Finally, the most relevant information along with original visuals obtained from the retrieval process will be used as the context for Qwen2.5-VL-7B-Instruct \cite{Bai2025Qwen25VLTR}, running locally on Ollama \cite{morgan2026ollama}, to generate output based on the research documents.

\cleardoublepage
\section{Project Requirements}
Our project aims to accomplish several objectives; therefore, a clear overview of the project requirements will be presented in this section. The requirements are divided into two main parts. First, the functional requirements specify the system functions and what it should do. Second, the non-functional requirements describe the system constraints and qualities. In addition, software and hardware requirements are discussed to specify the resources needed to develop the system.

\subsection{Functional Requirements}
\begin{enumerate}
  \item To support O1, the system shall allow users to upload and view research documents through the web-based interface.
  \item To support O3 and O4, the system shall allow users to ask questions about research documents using natural language.
  \item To support O2, the system shall process research documents locally without transferring document content to external models or tools running on external servers.
  \item To support O2, the system shall process both textual and visual content within research documents, including figures, tables, charts, and diagrams .
  \item To support O2, the system shall retrieve relevant textual and visual information independently from research documents based on users’ questions.
  \item To support O3 and O4, the system shall combine the relevant textual and visual information retrieved from research documents to support answer generation.
  \item To support O3, the system shall generate answers locally using a fine-tuned local language model. 
\end{enumerate}

\subsection{Non-Functional Requirements}
\begin{enumerate}
  \item Performance (O2, O4): The system shall perform document retrieval and respond to users' questions within 60 seconds for at least 80\% of test queries in the defined hardware environment.
  \item Security and Privacy (O1, O2): The system shall process and store research documents locally and shall not send document content, retrieved information, or responses to external servers or \ac{AI} services in the cloud.
  \item Usability (O1): The system shall offer an intuitive web-based interface, allowing users to upload, view, and manage research documents in order to achieve at least a 90\% success rate in performing tasks during usability testing.
  \item Availability (O1, O4): The system shall be available at least 95\% of the time during planned testing sessions, provided the local server and necessary hardware are operational.
  \item Retrieval Accuracy (O2): The system shall demonstrate at least 60\% retrieval accuracy for relevant textual and visual information from the research documents according to the selected evaluation dataset.
  \item Question Answering Accuracy (O3, O4): The system shall demonstrate at least 65\% question answering accuracy using the fine-tuned local model and the developed multimodal retrieval pipeline.
  \item Reliability (O1, O4): The system shall successfully complete at least 95\% of all valid requests for document processing and question answering without any unexpected system failures.
  \item Offline Operation (O2, O3): The system shall perform document processing, information retrieval, and answer generation without requiring an active internet connection after the necessary models and dependencies have been installed.
\end{enumerate}

\subsubsection{Hardware Requirements}
\begin{enumerate}
  \item Processor (O1, O2, O4): The system shall run on a computer featuring a modern multi-core processor, such as the Intel Core i7, AMD Ryzen 7, or a similar one, to ensure document processing, indexing, and general functioning of the system.
  \item Graphics Processing Unit (O2, O3, O4): The system shall be provided with a suitable \ac{GPU} with at least 8 to 12 GB of VRAM for fast visual document processing, multimodal inference, and fine-tuning. Depending on the chosen fine-tuning strategy, additional \ac{GPU} memory may be required.
  \item Random Access Memory (O2, O3, O4): The system shall have no less than 8 to 12 GB of RAM to allow for the parallel operation of document processing libraries, information retrieval modules, and locally stored \ac{AI} models.
  \item Storage (O1, O2, O3): The system shall have no less than 100 GB of SSD storage to store research documents, extracted data, vector databases, local \ac{AI} models, and associated files
\end{enumerate}

\subsubsection{Software Requirements}
\begin{enumerate}
  \item Programming Environment (O1-O4): The system shall use Python as the primary backend programming language to integrate document-processing libraries, retrieval techniques, and AI models.
  \item Document Processing Tools (O2): The system shall use MinerU for layout-aware document parsing, PyMuPDF for PDF processing, and Tesseract \ac{OCR} for extracting text from scanned research documents. These tools support the extraction of textual and visual document elements.
  \item Text Retrieval Tools (O2): The system shall use BGE-M3 for generating semantic embeddings, ChromaDB for storing and retrieving vectors, and \ac{BM25} for keyword retrieval, making it possible for the system to find relevant textual information.
  \item Visual Retrieval Model (O2): The system shall use ColQwen2 to generate visual representations and independently retrieve relevant figures, charts, and other visual information from research documents.
  \item Multimodal Retrieval Integration (O2, O4): The system shall implement the \ac{RRF} algorithm to combine and rank retrieved textual and visual information, improving the relevance of the evidence supplied to the generation model.
  \item Local \ac{VLM} (O3, O4): The system shall use Qwen2.5-VL-7B-Instruct to understand visual and textual evidence and generate answers based on retrieved research information. The model shall be adapted through fine-tuning to improve question-answering accuracy.
  \item Local Model Execution (O3, O4): The system shall use Ollama to execute the compatible local \ac{VLM} model for response generation without relying on external cloud-based \ac{VLM} services.
  \item Web Interface (O1, O4): The system shall use HTML, CSS, JavaScript, and PHP languages to implement a web interface that will enable users to engage with the research repository through the local backend and generate responses to their questions.
\end{enumerate}

\cleardoublepage
\section{Project Timeline}
This section represents the project milestones timeline. Each milestone is represented with a descriptive table and a Gantt chart. The separation of milestone timelines provides a clear overview of milestone scheduling and time distribution over tasks. Some factors have affected the selection of task durations. As a result, some tasks have a longer duration than others due to task complexity, exams, and assignments in other courses.

\begin{table}[H]
\centering
\begin{tabular}{|p{0.26\textwidth}| p{0.13\textwidth}| p{0.19\textwidth}| p{0.19\textwidth}| p{0.09\textwidth}|}
\hline
\textbf{Task} & \textbf{Assigned To} & \textbf{Start Date} & \textbf{End Date} & \textbf{Duration} \\
\hline
Milestone 1 & All Members & August 24, 2026 & October 29, 2026 & 67 days \\ 
\hline
Project Understanding & All Members & August 24, 2026 & September 5, 2026 & 10 days \\
\hline
Literature Review, Background, and Methodology & All Members & September 6, 2026 & September 30, 2026 & 25 days \\ 
\hline
Writing Introduction and Scope \& Deliverables & All Members & September 28, 2026 & October 7, 2026 & 10 days \\ 
\hline
Writing Project Requirements and Abstract & All Members & October 8, 2026 & October 12, 2026 & 5 days \\ 
\hline
Draft Submission to Supervisor and Review Feedback & All Members & October 13, 2026 & October 17, 2026 & 5 days \\ 
\hline
Milestone-1 Submission & All Members & October 18, 2026 & October 22, 2026 & 5 days \\ 
\hline
Milestone-1 Presentation & All Members & October 25, 2026 & October 29, 2026 & 5 days \\ 
\hline
\end{tabular}
\caption[Milestone 1 Task Schedule]{Milestone 1 Task Schedule.}
\label{tab:milestone1_schedule}
\end{table}

\begin{figure}[H]
\centering
\includegraphics[width=1\textwidth]{Milestone 1 Gantt chart.jpg}
\caption[Milestone 1 Gantt chart]{Milestone 1 Gantt chart.}
\label{fig:milestone1_gannt_chart}
\end{figure}

\begin{table}[H]
\centering
\begin{tabular}{|p{0.26\textwidth}| p{0.13\textwidth}| p{0.19\textwidth}| p{0.19\textwidth}| p{0.09\textwidth}|}
\hline
\textbf{Task} & \textbf{Assigned To} & \textbf{Start Date} & \textbf{End Date} & \textbf{Duration} \\ 
\hline
Milestone 2 & All Members & October 25, 2026 & December 17, 2026 & 54 days \\ 
\hline
System Architecture and Diagrams Design & All Members & October 25, 2026 & November 10, 2026 & 17 days \\ 
\hline
Comparison of Alternative Solutions & All Members & November 11, 2026 & November 16, 2026 & 6 days \\ 
\hline
Details of Tools and Programming Languages & All Members & November 17, 2026 & November 25, 2026 & 9 days \\ 
\hline
Initial Results & All Members & November 26, 2026 & November 30, 2026 & 5 days \\ 
\hline
Draft Submission to Supervisor and Review Feedback & All Members & December 1, 2026 & December 5, 2026 & 5 days \\ 
\hline
Milestone-2 Submission & All Members & December 6, 2026 & December 10, 2026 & 5 days \\ 
\hline
Milestone-2 Presentation & All Members & December 13, 2026 & December 17, 2026 & 5 days \\ 
\hline
\end{tabular}
\caption[Milestone 2 Task Schedule]{Milestone 2 Task Schedule.}
\label{tab:milestone2_schedule}
\end{table}

\begin{figure}[H]
\centering
\includegraphics[width=1\textwidth]{Milestone 2 Gantt chart.jpg}
\caption[Milestone 2 Gantt chart]{Milestone 2 Gantt chart.}
\label{fig:milestone2_gannt_chart}
\end{figure}

%% ----------------------------- References -----------------------------
\cleardoublepage
\printbibliography
\addcontentsline{toc}{section}{References}

%% ----------------------------- Appendices -----------------------------
% \appendix
% \cleardoublepage
%% ----------------------------- AppendixA -----------------------------
% \section*{Appendix A: Additional Materials}
% \addcontentsline{toc}{section}{Appendix A: Additional Materials}

\end{document}